\begin{abstract}
\noindent UREX is a handheld device that measures bus tyre tread depth by laser triangulation: a camera looks down at
the tread, a line laser strikes it at $35^\circ$, and the line shifts by about 2\,px for every millimetre of groove
depth. Our existing software works on good photos but reports a depth on most bad ones. We built a ray-traced
simulator of the printed holder and a benchmark with exact ground truth (five test sets, 800 frames), and used it,
together with 86 real hand-held photos, to compare ten families of algorithms: sub-pixel peak detectors,
dynamic-programming line tracking, a small U-Net, a learned candidate classifier, a robust depth stage, calibration
models, a quality gate, laser on/off subtraction, an end-to-end network, and a combination. The best single-frame
pipeline is almost entirely classical: Viterbi line tracking, a depth stage that fits the tread in millimetres and
clusters the groove floor, and two logistic-regression gates that can ask for a retake. Under simulated bad
conditions it reduces the mean depth error from 1.52 to 0.25\,mm, gross errors ($>1$\,mm) from 36\,\% to 2.7\,\%, while
answering 89\,\% of presses; across the three hardest sets it passes none of 41 illegal tyres where the baseline
passes 9; on real photos it stops reporting
depths on frames with no laser (cross-validated: 0 of 55) and agrees with the reference depths to about one pixel.
An independent audit reproduced every number exactly and found no training on test data. Networks that replace the geometry scored well on data like their training set and failed on other units and on
real photos. Laser on/off subtraction, which the firmware can already capture, would roughly halve the remaining error in simulation. A
preliminary hand-held test on a tyre of about 9\,mm gave repeatable pixel readings but millimetre estimates 2.4--4.2
times too large, because the camera--laser geometry was not fixed. Real-world accuracy is therefore still unmeasured;
calibrated measurements with the printed holder are the most valuable next step.
\end{abstract}
